\summarychapter{44}{New Creeds for Breakfast}{The Ecumenical Movement}

\begin{inbrief}
For four centuries the force in Protestantism ran outward, toward more and more denominations. In
the twentieth century a contrary force drew Christians toward cooperation, merger and united action.
The chapter traces this ecumenism from the Evangelical Alliance through the Edinburgh conference of
1910 to the World Council of Churches, reviews denominational mergers and confessional alliances, and
sets beside them the evangelical way to unity through common evangelism.
\question{How have Christians expressed the spirit of unity in the twentieth century?}
\end{inbrief}

Shortly before the 1961 assembly of the World Council of Churches at New Delhi, its general
secretary, Willem Adolph Visser 't Hooft, breakfasted in a Leningrad hotel with Russian Orthodox
leaders. The Council then defined itself as ``a fellowship of churches which accept our Lord Jesus
Christ as God and Saviour.'' The Russians objected that this ignored the Trinitarian basis of the
faith, and he recalled that Protestants had often missed any mention of Scripture. ``So I took the
breakfast menu and wrote out a new formula.'' A few months later the Council adopted it: a fellowship
of churches ``which confess the Lord Jesus Christ as God and Saviour according to the Holy Scriptures
and therefore seek to fulfill together their common calling to the glory of one God, Father, Son and
Holy Spirit.'' The move brought thirty million Russian Orthodox into the Council and changed the face
of the ecumenical movement.

In the sixteenth century the Reformation churches had four main divisions, Lutheran, Reformed,
Anabaptist and Anglican. Further denominations soon appeared, most of them founded by believers
convinced of the importance of some particular teaching of Scripture, and by the twentieth century
there were more than two hundred in the United States alone. The force was centrifugal, independent
and sometimes divisive. The new, centripetal force is called ecumenism.

\ssection{Movements toward Christian Unity}
\emph{Ecumenical} means worldwide or universal, and applied to the churches it implies the oneness
of Christians in the faith wherever they are. That unity may be understood as a spiritual reality
apart from human organizations, as Evangelicals tend to argue, or as something to be built by
federating churches or merging denominations. The spirit of unity is ecumenicity, the organizational
effort the ecumenical movement, and the creation of national and world councils of churches conciliar
ecumenism. Christians differ about doctrine, morality, worship and organization, and they hold their
views as convictions, not opinions. They therefore disagree about their divisions too: some defend
denominational distinctives, others call them a sin and a scandal.

The first significant modern effort at Protestant cooperation was the Evangelical Alliance, organized
in London in 1846 by believers from fifty evangelical bodies in England and America, with branches
later in nine European countries. It promoted religious liberty and joint activity, but its
enthusiasm cooled toward the end of the century. As this alliance of individuals flagged, a
federation of churches appeared. In 1908 thirty-one American denominations formed the Federal Council
of Churches, active in pronouncements on social, economic and political questions and criticized by
many conservatives for its liberal theology; in 1950 it was absorbed into the larger National Council
of Churches of Christ.

The most ambitious expression is the World Council of Churches, formed at Amsterdam in 1948. Shelley
pictures it as a great river fed by three tributaries: the International Missionary Council, the
Conference on Life and Work, and the Conference on Faith and Order. All three rise at the peak of
modern ecumenical history, the International Missionary Conference at Edinburgh in 1910, where more
than a thousand delegates met to consider world missions outside the Roman Catholic sphere and, in
discussing missions, discovered a great sense of unity. Four leaders tower over the early decades:
the American John R.~Mott, the Canadian Charles H.~Brent, the Swede Nathan S\"oderblom and the
Dutchman Visser 't Hooft.

Mott (1865--1955), a Methodist layman, joined deep faith and evangelistic zeal to a commanding power
over public assemblies. At twenty-three he became student secretary of the International Committee
of the YMCA. To coordinate student work he founded in Sweden the World's Student Christian
Federation, linking the student ministries of America, Britain, Scandinavia and Germany. The hundreds
of contacts this gave him made him the natural chairman at Edinburgh. Afterward he chaired the
Continuation Committee, and when the International Missionary Council was created in 1921 he was its
chairman for twenty years. No leader of the Council's movement did more for the spread of
Christianity.

Brent (1862--1929), a Canadian Anglican and missionary in the Philippines, was more aware than Mott
of the doctrinal differences that divide the churches, and he saw Anglicanism as a bridge across
them. When Edinburgh voted to seek ways of overcoming conflicts among the mission societies, he
carried the proposal to the triennial convention of his church at Chicago, which appointed a
committee to invite ``all Churches which accept Jesus Christ as God and Saviour'' to conferences on
``all questions pertaining to Faith and Order of the Church of Christ.'' Because of the First World
War and its aftermath the first World Conference on Faith and Order met only in August 1927, at
Lausanne, where 150 representatives of sixty-nine denominations passed resolutions that laid
foundations for the future World Council. Brent held that cooperation is possible only on agreement
in the essentials of the faith. Disunity is at bottom creedal, and until those differences are
resolved there will be no authentic unity. Faith and Order came to mean belief and worship.

\ssection{The Life and Work Movement}
S\"oderblom (1866--1931), Lutheran archbishop of Uppsala, founded and promoted the Life and Work
movement. When the king unexpectedly appointed him in 1914, conservatives questioned his orthodoxy,
and with reason. He rejected belief in the divine and human natures of Christ as unacceptable to
modern man. He regarded revelation as a continuing process not limited to the apostolic age, and
held that true religion lies in moral character, in what a person is and does, not in what he
believes. He therefore looked for unity through history, not through doctrinal agreement. Each
Christian group should respect the others and share its differences with them, and the one
revelation would gradually unfold through the ages. His real achievement, however, was in bringing
Christians together in common concern. He was the mind behind the first Conference on Life and Work
at Stockholm in August 1925, where five hundred delegates from thirty-nine countries and ninety-one
denominations concluded that the problems of social morality are too great for individual effort and
that the Christian community must take responsibility for the common good.

By 1937 both conferences saw that the quest for unity needed a new and more inclusive body, and
meeting in Britain that year they jointly called for a World Council of Churches. Hitler and the war
delayed it. The first assembly met at Amsterdam in 1948 with 351 delegates of 147 denominations from
forty-four countries. The chief absentees were the Roman Catholics, many conservative Evangelicals
and the Russian Orthodox. The general secretary in these formative years was Visser 't Hooft
(1900--1985). He had followed in Mott's footsteps as secretary of the World's Committee of the YMCA
and then of the World's Student Christian Federation, and in 1938 he was the natural choice to lead
the Provisional Committee. Karl Barth influenced him deeply. Barth, he said, felt that the church had
nearly lost its soul in adjusting to historical trends and ``called the church again to be itself.''
The unofficial slogan of the founders was ``Let the Church be the Church.'' That did not mean running
away from the world. It meant that the church is more than an echo of the world's trends.

One of his favorite projects after the war was an Ecumenical Institute in Switzerland for training
leaders of the unity movement. At a dinner in the United States he described the plan to John
D.~Rockefeller Jr., who told him, ``You must ask for more money,'' and later gave about a million
dollars to establish the institute at Bossey. Thanks to his diplomacy an assembly of the Council
became a mosaic of cultures and concerns: theological conservatives from the Orthodox lands of
eastern Europe and the Middle East, secular theologians from Europe and North America, Evangelicals
from Europe and the Third World, confessional Lutherans from Scandinavia and liberation spokesmen
from Latin America. The Council did not claim to be a super-church. Under its constitution it cannot
legislate for its members, and its aim is understanding, cooperation and unity wherever possible, so
that its statements are inevitably compromises, as the Leningrad breakfast shows.

Later assemblies met at Evanston, Illinois (1954), New Delhi (1961), Uppsala (1968) and Nairobi
(1975). At New Delhi the Russian Orthodox Church joined and the International Missionary Council came
under the World Council's administration. With the years the doctrinal emphasis of Faith and Order
declined and the social concerns of Life and Work grew. Uppsala and Nairobi reflected the inroads of
secularism in the churches. Racism, war and peace, poverty and unemployment, alcoholism and drug
addiction and the women's liberation movement so occupied the councils of churches that the specific
task of reuniting separated church bodies was left more and more to other agencies.

\ssection{Denominational Mergers}
Zeal for unity also produced mergers and international alliances. In the United States, the home of
denominationalism, more than thirty mergers were completed between 1900 and 1970, among them the
United Methodists and the United Presbyterians. Some reunited bodies that had split in the
nineteenth century over issues such as race (the Methodists), and others joined groups of one
tradition that had come from different countries (the Lutherans). Abroad the most significant was
the Church of South India, formed in 1947 from the Anglican Church of India, Burma and Ceylon, the
South India Province of the Methodist Church and the South India United Church, itself a union of
Presbyterian, Congregational and Dutch Reformed bodies.

A plan of similar size would have created an American church of nineteen million members. In 1960
Eugene Carson Blake, chief executive of the United Presbyterian Church (the Northern Presbyterians)
and later general secretary of the World Council, proposed that his church and the Protestant
Episcopal Church together invite the Methodists and the United Church of Christ to form a new
church. ``I don't believe it is God's will,'' he said, ``to have so many churches in the United
States.'' The Episcopal bishop James A.~Pike welcomed the idea: the Holy Ghost is on our side
whenever we break through the barriers between Christian bodies. At the Episcopal General Convention
in Detroit in September 1961 a whole day had been set aside for the bishops' debate and a crowd
packed the hall to watch the fireworks; the motion passed unanimously in about thirty seconds. ``I'm
quite speechless!'' said the presiding bishop, Arthur C.~Lichtenberger, who had prayed just before
the vote, ``Look mercifully, O Lord, on the broken body of Thy Church.'' The convention closed with
the bishops' appeal to Christians to work and pray without ceasing until the divisions that dishonor
the one Lord are done away. The Consultation on Church Union (COCU) then began to seek a plan,
later renaming itself the Church of Christ Uniting so that other denominations could join. Time
dampened the enthusiasm, and after twenty years the churches were still looking for the right
formula.

The century also brought world confessional alliances, in which churches of similar confession and
organization meet periodically for discussion and fellowship: the International Congregational
Council, the Mennonite World Conference, the World Methodist Conference, the Baptist World Alliance,
the Lutheran World Federation and the World Alliance of Reformed and Presbyterian Churches.

The most persistent critics of conciliar ecumenism have been the conservative Evangelicals. They know
that Jesus prayed for his disciples to be one, but they question unity by federation. They find the
World Council's doctrinal basis inadequate and doubt its commitment to evangelism, and they are
troubled by its growing involvement in the politics of developing nations. Since Uppsala the Council
has seemed to treat the unity of the church as a sign of the unity of mankind, which Evangelicals
think leads easily to humanistic goals for society in place of a distinctively Christian witness.
The difference between church and world becomes less one of faith and unbelief than of oppressed and
oppressor, and salvation becomes liberation.

\ssection{Evangelical Expressions of Unity}
Rooted in revivalism, Evangelicals have always stressed personal religious experience. They care
little for the proper form of the church and much for its mission. Whitefield spoke for them all
when, preaching from the courthouse balcony in Philadelphia, he asked Father Abraham whom he had in
heaven: Episcopalians, Presbyterians, Independents, Methodists? No; ``we don't know those names
here. All who are here are Christians.'' Then, he cried, ``God help us to forget party names and to
become Christians in deed and truth.''

In the early 1940s American Evangelicals created the National Association of Evangelicals and the
American Council of Christian Churches. Both were loyal to orthodox Christianity, but they differed
in structure and in their attitude to conciliar ecumenism. The American Council sharply criticized
the National and World Councils and everyone associated with them in any way. Internationally
Evangelicals turned to congresses that encouraged united effort in evangelism. Inspired by Billy
Graham's crusades, the World Congress on Evangelism at Berlin in 1966 drew participants from more
than a hundred countries and stimulated regional and national conferences on most continents. The
International Congress on World Evangelization at Lausanne in 1974 showed a new maturity. A group of
142 leaders under Graham's honorary chairmanship invited 2,700 participants to stimulate regional
groups and to frame the Lausanne Covenant, which the great majority signed.

The Covenant affirms ``that the church's visible unity in truth is God's purpose,'' for two reasons.
The theological one is that unity is a gift of God through the Spirit, made possible by the cross of
Christ, who ``is our peace.'' It may take many forms, and organic unions do not always preserve unity
``in truth.'' The practical one is that ``evangelism summons us to unity,'' since a gospel of
reconciliation cannot well be preached by the unreconciled. The congress created a forty-eight-member
Continuation Committee for World Evangelization, which chose an African, Gottfried Osei-Mensah, as
executive secretary and set out to encourage regional and national committees for evangelization in
every area. So, as the 1970s closed, the ecumenical spirit in the World Council had apparently turned
to social concerns, sometimes with openly political instruments, as the chief expression of unity,
while conservative Evangelicals aimed to restore evangelism to the center of the church's mission in
the hope that unity would follow.

\begin{keyfacts}
\keypeople{Willem A.~Visser 't Hooft; John R.~Mott; Charles H.~Brent; Nathan S\"oderblom; Karl Barth;
Eugene Carson Blake; James A.~Pike; Arthur C.~Lichtenberger; Billy Graham; Gottfried Osei-Mensah.}
\keydates{%
  \kd{1846}{Evangelical Alliance organized in London}
  \kd{1908}{Federal Council of Churches formed in the United States}
  \kd{1910}{International Missionary Conference at Edinburgh}
  \kd{1925}{First Conference on Life and Work, Stockholm}
  \kd{1927}{First World Conference on Faith and Order, Lausanne}
  \kd{1947}{Church of South India formed}
  \kd{1948}{World Council of Churches formed at Amsterdam}
  \kd{1961}{New Delhi assembly: new basis; Russian Orthodox Church joins}
  \kd{1966}{World Congress on Evangelism, Berlin}
  \kd{1974}{International Congress on World Evangelization and the Lausanne Covenant}}
\keyterms{Ecumenical; ecumenicity; ecumenical movement; conciliar ecumenism; Faith and Order; Life and
Work; International Missionary Council; World Council of Churches; denominational merger;
confessional alliances; Lausanne Covenant.}
\end{keyfacts}
